% from ega0-22-fr.tex
\oldpage[0\textsubscript{IV-1}]{190}

et l'hypothèse que $u_\alpha$ est injectif pour
$\alpha\geqslant\alpha_0$ implique que $\operatorname{Ker}(u)$ est contenu
dans l'\emph{intersection} des $\operatorname{Ker}(v_\alpha)$. Mais on a vu
(21.8.3) que cette intersection est nulle en vertu de l'hypothèse
$\bigcap_\alpha k_\alpha(k^p)=k^p$, ce qui achève la démonstration.
\end{proof}

\begin{proposition}[22.2.9]
\label{0.22.2.9-fr}
Sous les hypothèses de (22.2.1), les conditions suivantes sont
équivalentes~:
\begin{enumerate}
  \item[a)] $A/\mathfrak m^2$ est une $k$-extension $k$-triviale de $K$ par
  $\mathfrak m/\mathfrak m^2$.
  \item[b)] $\chi_{A/k}=0$.
  \item[c)] L'homomorphisme canonique (20.5.11.2)
  \[
    \delta_{K/A/k}:\mathfrak m/\mathfrak m^2
    \longrightarrow\Omega^1_{A/k}\otimes_AK
  \]
  est injectif (en d'autres termes, la suite
  \begin{equation}
  \tag{22.2.9.1}
  \label{0.22.2.9.1-fr}
    0\longrightarrow\mathfrak m/\mathfrak m^2
    \longrightarrow\Omega^1_{A/k}\otimes_AK
    \longrightarrow\Omega^1_{K/k}\longrightarrow0
  \end{equation}
  est exacte).
\end{enumerate}
\end{proposition}

\begin{proof}
L'équivalence de $b)$ et $c)$ résulte de la suite exacte (20.6.22.1)~;
l'équivalence de $c)$ et $a)$ résulte de (20.5.12), la suite (22.2.9.1)
étant scindée si elle est exacte, puisqu'il s'agit d'espaces vectoriels.

On notera que si $K$ est \emph{séparable sur $k$}, les conditions
équivalentes de (22.2.9) sont satisfaites, en vertu de la définition de
$\chi_{A/k}$ (22.2.1) et de (20.6.3).
\end{proof}

\begin{corollary}[22.2.10]
\label{0.22.2.10-fr}
Sous les hypothèses de (22.2.1), considérons un sous-corps $k_0$ de $k$. Les
conditions suivantes sont équivalentes~:
\begin{enumerate}
  \item[a)] $A/\mathfrak m^2$ est une $k_0$-extension $k_0$-triviale de $K$
  par $\mathfrak m/\mathfrak m^2$.
  \item[b)] L'homomorphisme composé
  \[
    \Upsilon_{K/k_0}\longrightarrow\Upsilon_{K/k}
    \xrightarrow{\chi_{A/k}}\mathfrak m/\mathfrak m^2
  \]
  est nul.
  \item[c)] L'homomorphisme caractéristique
  $\chi_{A/k}:\Upsilon_{K/k}\to\mathfrak m/\mathfrak m^2$ se factorise en
  \[
    \Upsilon_{K/k}\longrightarrow\Upsilon_{K/k/k_0}
    \xrightarrow{\chi'}\mathfrak m/\mathfrak m^2.
  \]
\end{enumerate}
En outre, lorsque ces conditions sont remplies, l'homomorphisme $\chi'$ est
déterminé de façon unique et coïncide avec l'homomorphisme caractéristique
de la $k_0$-extension $k_0$-triviale $A/\mathfrak m^2$ de $K$ par
$\mathfrak m/\mathfrak m^2$.
\end{corollary}

\begin{proof}
Comme le composé dans $b)$ n'est autre que $\chi_{A/k_0}$, l'équivalence de
$a)$ et $b)$ résulte de (22.2.9)~; d'autre part, $b)$ et $c)$ sont
équivalentes en vertu de la suite exacte (20.6.18.1)
\[
  \Upsilon_{K/k_0}\longrightarrow\Upsilon_{K/k}
  \longrightarrow\Upsilon_{K/k/k_0}\longrightarrow0,
\]
et comme l'homomorphisme canonique
$\Upsilon_{K/k}\to\Upsilon_{K/k/k_0}$ est surjectif, cela implique l'unicité
de $\chi'$~; le fait que $\chi'$ est l'homomorphisme caractéristique de la
$k_0$-extension $A/\mathfrak m^2$ résulte de (20.6.22.2).
\end{proof}

\begin{corollary}[22.2.11]
\label{0.22.2.11-fr}
Sous les hypothèses de (22.2.1), soit $p$ l'exposant caractéristique de $k$,
et soit $(k_\alpha)$ une famille filtrante décroissante de sous-corps de $k$
telle que $\bigcap_\alpha k_\alpha(k^p)=k^p$. Si $\Upsilon_{K/k}$ est de rang
fini sur $K$, il existe un indice $\alpha$ tel que $A/\mathfrak m^2$ soit une
$k_\alpha$-extension $k_\alpha$-triviale de $K$ par
$\mathfrak m/\mathfrak m^2$.
\end{corollary}

\begin{proof}
On sait en effet (21.8.3) qu'il existe alors un indice $\alpha$ tel que
l'homomorphisme canonique
$\Upsilon_{K/k}\to\Upsilon_{K/k/k_\alpha}$ soit bijectif, donc la condition
$b)$ de (22.2.10) (où $k_\alpha$ remplace $k_0$) est satisfaite.
\end{proof}
